% !TeX root = ../../main.tex
\section{Methodology}
\label{sec:methodology}

% Draft from the outline: Methodology. Subsection names follow the outline.

\subsection{Framework Overview}
\label{sec:framework}

Evaluating the candidate metrics requires real and simulated event streams acquired under comparable conditions. Because many event-camera simulators take images or video as input, we need paired recordings in which an event camera and a frame camera observe the same stimulus, with spatial and temporal alignment established through calibration. The frame recording then provides the input to the simulators, so that the simulated events can be compared directly with the real events after careful calibration. We also need repeated observations under nominally identical conditions, which show how much real recordings differ from one another and provide a floor for referencing and interpreting the differences between real and simulated events.

We therefore built a data collection and calibration framework around a controllable periodic stimulus. Real and simulated events are obtained in pairs from the same stimulus and converted into a common point-set representation, in which polarities are pooled, timestamps are scaled into the same units as pixel coordinates, and the stream is divided into short windows (Section~\ref{sec:representation}). MMD and SWD, which compare events at the level of distributions, are computed on this representation, together with Chamfer distance as an independent verification. Before these metrics are used to compare real and simulated events, their responses are examined with controlled perturbations of the real recordings.

The first stage records event data and frame video that are spatially aligned and temporally synchronized. The event camera is a Prophesee EVK4 with a resolution of 1280 by 720 pixels, and the frame camera is an Allied Vision Mako U-130b, which records the video at 30 frames per second. The two cameras are mounted side by side and fitted with identical 8~mm lenses set to the same f-number. Because they view the scene from slightly different positions, their recordings are aligned in two steps, described below: spatial alignment and temporal synchronization.

% Leave a space for 

For data collection, the stimulus is an optical chopper, chosen because it produces a periodic signal whose frequency can be adjusted. Each recording contains many nominally identical rotations, which provide repeated samples of the same stimulus for estimating the noise floor, and changing the rotation speed gives a single physically controlled variable along which the metrics can be tested. We record five speeds, named F1 to F5 after their aperture frequencies of 1 to 5~Hz, which correspond to rotation periods from about 1.96~s to 0.40~s. The chopper motor runs freely rather than locked to the cameras, and its speed can fluctuate, so the rotation period of each recording is measured from the real events instead of being taken from the set frequency. The wheel has two apertures and therefore four blade edges per rotation, and the event rate peaks at each edge, so the period is obtained from the strongest frequency of the event rate. We selected an optical chopper because it rotates relatively stably at the low speeds our experiments require, since faster rotation would cause serious motion blur in the frame camera. The fan we considered was less stable in this range, while a freely swinging pendulum loses amplitude over time and its speed varies throughout each cycle.

\subsubsection{Spatial Alignment}
\label{sec:spatial-alignment}

The two cameras are spatially aligned in two steps, first by a homography estimated from a checkerboard that maps the frame-camera image onto the event-camera image plane, and then by a refinement based on the chopper disk's geometry. For the first step, each camera is calibrated with a printed checkerboard to estimate its intrinsic parameters and lens distortion~\cite{zhang2000flexible}. The event camera does not record images, so its checkerboard recording is first converted into intensity frames with E2VID~\cite{rebecq2021e2vid}, after which the same corner detection and calibration can be applied. The reprojection errors are 0.21 pixels for the frame camera and 0.81 pixels for the event camera. The homography is then computed from the checkerboard corners detected in one undistorted image pair. Every frame of the video is undistorted and warped through this homography into the 1280 by 720 event-camera image plane, and these warped frames are the input to both simulators. Lens distortions are then corrected for the real events with the event camera's intrinsic parameters, so that the real and simulated events share the same pixel coordinates.

The homography alone does not align the chopper precisely, because it maps points correctly only on the plane in which it was estimated. First, it is difficult to place the checkerboard exactly at the depth of the chopper disk during calibration, so points on the disk are mapped with a small residual offset. Second, because the two cameras view the scene from slightly different positions, any part of the scene at a different depth is also affected by parallax, which a single homography cannot correct. After the homography, the chopper disk fitted in the real events, which is about 426 pixels in diameter, is still offset from the disk in the warped video by 20 to 22 pixels horizontally and 1 to 4 pixels vertically, about 5 percent of the disk area and under 1 percent of its diameter, and differs in size by 1 to 2 percent. This residual is nearly the same in all five recordings, which means that it comes from the experiment setup rather than from an individual recording.

Therefore, a second step is needed to refine the alignment using the chopper disk itself, whose circular outline appears as an ellipse in both cameras. The disk is located by accumulating the events of a full rotation into an image, separately for the real events and for the events that v2e simulates from the warped video, which share the video's coordinates by construction. A morphological opening removes thin structures attached to the disk, such as cables and mounts, and an ellipse is fitted to the convex hull of the remaining region. Because the disk is viewed nearly head-on, both fitted ellipses are almost circular, with axis ratios below 1.02, so their orientations are not well defined. The correction therefore consists only of a shift and a uniform scale that match the center and mean diameter of the real disk to those of the simulated one, and it is applied to all real events. Since this correction is estimated on the chopper disk, the second step is valid only near the plane of the disk rather than throughout the scene, so both the real and the simulated events are restricted to the disk region before they are compared. We will only compare the events from this disk region. 

\subsubsection{Temporal Synchronization}
\label{sec:temporal-sync}

The two recordings are also synchronized in time in two steps. The first step involves shared external trigger pulses and the second step is a refinement based on the periodic motion frequency of the chopper. External triggers are generated from an Analog Discovery 2 every frame. Each pulse triggers a frame captures in the conventional frame camera and is recorded by the event camera as a timestamped trigger event. The real event stream is then cut into intervals with these external triggers, from the first trigger to the end of the frame interval, until next trigger event. So the real event stream is processed to the same duration of the corresponding video recording, which is 15 seconds for each trial. However, the two cameras do not respond to a trigger with the same processing delay, so an offset remains between the two recordings which may impact the temporal synchronization result.   

The remaining offset is measured from the chopper's motion itself, which appears in both real and v2e events. Within the fitted disk's mask, each event's position is converted into a polar angle around the disk center, and the edges' angles are averaged in 2.5ms bins. This gives a sequence of angle signals for the chopper blade in rotation of each event stream. Because ON and OFF events are generated slightly differently for each event pixel and they mark opposite edges of a blade, they are processed separately. If the two angle signals different only by an offset, that offset can be interpreted as a phase difference between the two periodic signals in frequency space. The offset can therefore be calculated from the phase difference. A phase, however, defines the delay which is smaller than a whole period in value, so if the offset is larger than a full period, it will only report the remainder after the full offset is divided by a period duration. The results from the cross correlation of the two signals can be used to find if the difference is larger than a full period. In all five recordings, it selected the closest cycle, so the offset is limited to the phase offset itself. The offset is estimated separately from ON and OFF events, which agree to within 1.5 to 5ms, and their average is subtracted from all real-event timestamps.  


\subsection{Event Representation}
\label{sec:representation}

To measure the difference between two event streams, the events need a representation in which a distance can be computed without dropping the information they carry. Each event is recorded as $(x, y, t, p)$, with pixel coordinates, a timestamp in microseconds, and a binary polarity, so its components differ in type and unit. Many representations have been proposed for processing events, such as event frames, voxel grids, and time surfaces, but each is usually chosen as the input format for a particular algorithm and discards part of this information~\cite{gallego2022event}. Event frames and voxel grids quantize the timestamps into time bins, and time surfaces keep only the most recent event at each pixel. We therefore keep the events as points in space and time, $(x, y, t)$, which preserves every event with its exact position and timestamp~\cite{gallego2022event}. Polarity is not included, so ON and OFF events are pooled into a single point set. The event streams are divided into short temporal windows, each forming one point set. The window length is a parameter that can be chosen for the application (Section~\ref{sec:parameters}). Event-based eye tracking methods have grouped events over windows from about 1 to 40~ms~\cite{bonazzi2024retina,groenen2025gazescrnn,wang2024eyetracking,zhao2023eveye}, and in this study we use 5~ms. We use windows for two reasons. First, real applications usually process event data in small batches rather than as a whole stream~\cite{gallego2022event}, so a window also keeps each distance local in time. Second, windows keep the computation feasible, because metrics such as MMD compare every pair of points, so their cost grows with the square of the number of events, making it very computationally expensive to calculate over the whole dataset.

In these points, however, spatial coordinates are measured in discrete pixels positions bounded by the sensor resolution, while timestamps are in microseconds which is continuous and their range grows with the length of the recording. Therefore, after calibration, the spatial coordinates and timestamps of the events still need to be converted into a unified space to calculate interpretable distances. Otherwise, temporal units will dominate the distance, making the result always biased to the temporal difference. For example, within a 5 millisecond window, the timestamps varies in 5000 units, but the optical chopper is only about 426 pixels in diameter, so a time difference of about 400 microsecond will be weighted as much as a displace across the whole disk. 

We considered two ways to set this conversion. The first normalizes each axis globally, for example with a z-score or by dividing it by its range. The second relates time to space through the speed of the stimulus, either measured from the data or set to a typical value for a type of stimulus. We use the second approach with a empirical value to make an estimation, because both a global normalization and a speed measured from the data vary from one recording to another, whereas a fixed ratio keeps distances comparable across recordings and corresponds to a physical speed. Since our target application is eye tracking, this speed is taken from a possible large saccade. The fastest eye movement. Saccade duration grows with amplitude as $D = 2.2A +21$~ms~ \cite{bahill1975mainsequence,carpenter1988movements,gibaldi2021saccade}, so a $15^\circ$ saccade lasts about 54~ms. If such a saccade moves the image across the 1280-pixel width of the sensor, it travels one pixel in about 42~\textmu s, and we therefore scale the timestamps by $s = 1/42$ pixel per microsecond. Each event is therefore represented by the point $(x, y, st)$.


\subsection{Maximum Mean Discrepancy}
\label{sec:mmd}

% --- To be edited myself ---

MMD compares two event distributions through the average kernel similarity within and between their point sets~\cite{gretton2012kernel}. For two windows with scaled point sets $X=\{u_i\}_{i=1}^{m}$ and $Y=\{v_j\}_{j=1}^{n}$, with every event weighted equally, we use the unbiased estimator
\begin{equation}
\begin{split}
  \widehat{\mathrm{MMD}}_u^{\,2}(X,Y) = {}& \frac{1}{m(m-1)} \sum_{i \ne i'} k(u_i,u_{i'}) \\
  &+ \frac{1}{n(n-1)} \sum_{j \ne j'} k(v_j,v_{j'}) \\
  &- \frac{2}{mn} \sum_{i=1}^{m} \sum_{j=1}^{n} k(u_i,v_j),
\end{split}
\label{eq:mmd}
\end{equation}

where $k$ is the kernel that measures the similarity between two points. The within-set sums exclude each event's comparison with itself. Including it would add about $1/m + 1/n$ to the estimate, so two sets drawn from the same distribution would appear to differ according to their numbers of events. The unbiased estimate can be negative when the two distributions are close. We report the distance $\sqrt{\max(\widehat{\mathrm{MMD}}_u^{\,2},0)}$, and keep the signed estimate for statistics over many comparisons, where setting negative values to zero would bias the result upward.

The similarity between points is measured with a Gaussian kernel, whose width sets the scale at which differences between events contribute to MMD. The kernel is $k(u,v)=\exp\!\left(-\|u-v\|^2/(2\sigma^2)\right)$. Instead of setting $\sigma$ directly, we specify the distance $d$ at which the similarity falls to one half, which gives $\sigma = d/\sqrt{2\ln 2} \approx 0.85\,d$. Because the timestamps are scaled by $s$, $d$ applies to space and time alike: with $d = 15$~pixels, a displacement of 15 pixels and a time difference of 630~\textmu s both reduce the similarity to one half. We use three kernels, with $d = 3$, 15, and 75~pixels ($\sigma \approx 2.5$, 12.7, and 63.7~pixels), to compare local detail with broader event structure. Their choice is discussed in Section~\ref{sec:parameters}.

\subsection{Sliced Wasserstein Distance}
\label{sec:swd}

% --- To be edited myself ---

SWD compares two event distributions through the cost of transporting one onto the other along many one-dimensional projections~\cite{bonneel2015sliced}. Each point set is treated as a distribution with equal weight per event, and both are projected onto the same $L$ unit directions $\theta_l$ in the scaled space--time. Along each direction, the Wasserstein distance between the projected values is computed from their sorted order, which is far cheaper than transport in the original three dimensions. With order $p = 2$,
\begin{equation}
  \mathrm{SWD}_2(X,Y) = \left( \frac{1}{L} \sum_{l=1}^{L} W_2^2\!\left(\theta_l^{\top}X,\ \theta_l^{\top}Y\right) \right)^{1/2},
  \label{eq:swd}
\end{equation}
where $\theta_l^{\top}X$ denotes the projected values of the points in $X$. We use $L = 100$ directions drawn at random with a fixed seed, as implemented in the Python Optimal Transport library~\cite{flamary2021pot}, so every comparison uses the same directions. A finite set of directions only approximates SWD, and the fixed seed makes the result reproducible without removing this approximation error. Unlike MMD, SWD has no scale parameter, and its value is a distance in the scaled space, measured in pixels.

\subsection{Chamfer Distance}
\label{sec:chamfer}

% --- To be edited myself ---

Chamfer distance compares the two point sets through nearest neighbors~\cite{fan2017pointset}. We apply it to the same scaled point sets as MMD and SWD, and use the sum of the mean nearest-neighbor Euclidean distances in both directions,
\begin{equation}
  \mathrm{CD}(X,Y) = \frac{1}{m}\sum_{i=1}^{m} \min_{j} \|u_i - v_j\| + \frac{1}{n}\sum_{j=1}^{n} \min_{i} \|v_j - u_i\|.
  \label{eq:chamfer}
\end{equation}
The distances are not squared, and the sum is not divided by two. This is the same form that PECS used to compare real and simulated events, there on four-dimensional points that also include polarity~\cite{han2024pecs}. Like SWD, its value is a distance in the scaled space, measured in pixels. As discussed in Section~\ref{sec:background}, its value also changes with the number of events, so we use it as an independent verification rather than as one of the main metrics.

\subsection{Choice of Parameters}
\label{sec:parameters}

\subsection{Metric Evaluation Procedure}
\label{sec:procedure}
